\section{Thermal stability of a retuned zero}
\label{sec:thermal}

\subsection{An operator zero and a ground-state zero}

Temperature tests whether exchange retuning suppresses all low-energy
transitions or only those accessible from the ground state.
We keep $\tau=T/J_{\rm eff}$ fixed as
$\beta\to0$, where $J_{\rm eff}=\beta c_1(\alpha_*)^2$ and $k_B=1$.
Thus the thermal ensemble remains inside the local ground multiplets.
Put $h=\sum_i\bm T_i\cdot\bm T_{i+1}$ and let $\epsilon_i$ be its
dimensionless eigenvalues. For the Hermitian leading $Q=\pi$ probe,
write $D=D_\pi$, $K=K_q$, and define
\begin{align}
 R_\kappa&=K+\kappa A'D,\notag\\
 C_\kappa(\tau)&=\sum_{\epsilon_f>\epsilon_i}
 p_i\left|\langle f|R_\kappa|i\rangle\right|^2,
 \qquad p_i=\frac{e^{-\epsilon_i/\tau}}{Z}.
 \label{eq:thermal-weight}
\end{align}
This coefficient describes transitions between distinct effective energy
levels, with frequencies on the $J_{\rm eff}$ scale.
At fixed finite size, choose $\eta>0$ smaller than
the smallest nonzero level spacing of $h$, and let $W_{+,\eta}$ include
only low-band transitions with $\omega>\eta J_{\rm eff}$. Then
$W_{+,\eta}=\beta^2C_\kappa+O(\beta^3)$.
We call this the regular positive-frequency response.
Higher-order splittings within an equal-energy block instead produce
quasielastic lines at $\omega=O(\beta^2)$, outside this leading coefficient.
The minimum $C_{\min}(\tau)$ allows a separate real retuning at each
temperature, including any temperature-dependent displacement of the zero.

At any $\tau>0$ every state of a finite effective chain has positive
Boltzmann weight. Nonnegativity of the summands and Hermiticity imply
\begin{equation}
 C_\kappa(\tau)=0\quad\Longleftrightarrow\quad[h,R_\kappa]=0.
 \label{eq:thermal-criterion}
\end{equation}
Thus complete thermal cancellation requires the retuned probe to act
only within equal-energy blocks. Equation~\eqref{eq:thermal-criterion}
reduces the thermal question to an operator calculation.
For a doublet, the stronger identity
$K=\gamma D$ already proved above cancels $R_\kappa$ throughout the
retained space. Consequently $W_+=O(\beta^4)$ or smaller at every fixed
$\tau$ under the same periodic and reflection assumptions.

\subsection{A representation-dependent obstruction}

For the quartet root at $q=\pi$, define the unnormalized alternating spin
$\widetilde D=\sum_i(-1)^iT_i^z$ and
$Y=\sum_i(-1)^iX_{i,i+1}$, using Eq.~\eqref{eq:bond-identity}. Then
\begin{equation}
 R_\kappa=\frac{bY+(-2a+\kappa A')\widetilde D}{\sqrt{5N}}.
 \label{eq:thermal-source}
\end{equation}
Both $\widetilde D$ and $Y$ are spin vectors with momentum $\pi$,
odd under time reversal and even under the site-permutation reflection
$i\mapsto-i$. For $Y$, reflection reverses the bond orientation and
shifts its alternating sign, leaving their product unchanged.
Their different bond-channel amplitudes, rather than different symmetry
labels, control the cancellation. To test whether one dipole coefficient
can remove every inelastic matrix element, we evaluate the squared
commutator using a spectral moment~\cite{mueller1982sumrules}.
Write $u=-2a+\kappa A'$ in Eq.~\eqref{eq:thermal-source}.
For every finite even quartet chain with $N\ge4$, the uniformly traced,
two-sided second frequency moment is
\begin{align}
 M_2(u)&=\frac{1}{4^N}\operatorname{Tr}[h,R]^\dagger[h,R]\notag\\
 &=\frac{197}{16}b^2+\frac52(u+b)^2.
 \label{eq:quartet-moment}
\end{align}
The positive first term survives every choice of $u$.
Together with Eq.~\eqref{eq:thermal-criterion}, it proves
$C_{\min}(\tau)>0$ at every $\tau>0$ for these finite chains.
Even a ground state that can be made dark therefore acquires an
irreducible leading thermal response.

The proof uses orthogonality of local operator supports, so the same
minimum applies at every finite even $N\ge4$.
Appendix~\ref{app:general-moment} derives the spin-$s$ expression and
shows how its minimum vanishes for a doublet.
Here the uniform trace covers the retained effective space, corresponding
to $J_{\rm eff}\ll T\ll\Dloc$ in the fixed-size weak-coupling limit.
Combining its second and fourth moments also bounds a finite-frequency
part of the quartet weight away from zero, uniformly in even $N$
(Appendix~\ref{app:weight-bound}). This latter result takes the effective
uniform-trace limit before increasing the chain length.

\subsection{The accessible local-control space}

Independent exchange controls within each cell leave this obstruction
intact. At the quartet root $q=\pi$, any bounded
SU(2)-invariant perturbation of order $\beta$ confined to an individual
cell changes its projected vector probe only by a multiple of $T_i^z$.
The isolated cell contains a single spin irrep; its projected scalar
perturbation shifts all components equally. The leading uniform $h$ and
the coefficient $b$ are therefore unchanged. Even independent local
controls in every cell can only enlarge the leading probe to
\begin{equation}
 R_{\{u_i\}}=\frac{bY+\sum_i(-1)^iu_iT_i^z}{\sqrt{\ell N}},
 \qquad u_i\in\mathbb R.
 \label{eq:local-control-probe}
\end{equation}
Here $\ell=5$ fixes the normalization inherited from the microscopic probe.
These controls preserve the local irrep and the leading intercell
Hamiltonian; their changes to either $h$ or $b$ enter at higher order.

Translation symmetry separates the $Q=\pi$ component, whose dipole
coefficient is $\bar u=N^{-1}\sum_i u_i$, from all other operator momenta.
Their thermal spectral weights add without interference. For either the
positive-frequency weight or any fixed positive-frequency window,
\begin{equation}
 F(R_{\{u_i\}})=F(R_{\bar u})+
 F(R_{\{u_i\}}-R_{\bar u})\ge F(R_{\bar u}).
 \label{eq:local-control-projection}
\end{equation}
Here $F$ uses a Gibbs state of the uniform leading $h$, or its uniform
trace. The minimum over all real $u_i$ therefore equals the uniform
minimum. Spatial variation adds spectral weight at other momenta but
cannot compensate the residual $Q=\pi$ response.
Appendix~\ref{app:local-controls} gives the full proof.

\section{Thermally activated response and microscopic spectra}

\subsection{Exact activation from a dark four-cell ground state}

The four-cell quartet ring makes the contrast explicit: the retuned
leading probe annihilates its ground state, yet thermally occupied
states scatter. Write $\bm A=\bm T_0+\bm T_2$ and
$\bm B=\bm T_1+\bm T_3$, so that $h=\bm A\cdot\bm B$.
Its ground state has $(j_A,j_B,J)=(3,3,0)$ and energy $-12$;
the first excited multiplet is $(3,3,1)$ at $-11$.
The zero-temperature optimum is $\lambda_0=-2a-3b$, with
$\lambda=-\kappa A'$. The resulting source
$b(Y+3\widetilde D)/\sqrt{20}$ annihilates the ground state exactly,
although it does not commute with $h$.

The first thermally occupied triplet opens two groups of transitions.
Their final sectors are $(3,3,2)$ at $\omega/J_{\rm eff}=2$ and
$(3,1,2),(1,3,2)$ at $\omega/J_{\rm eff}=7$.
Summed over the three initial magnetic components, the squared matrix
elements divided by $b^2$ are respectively $27/50$, $6/25$, and $6/25$.
Consequently,
\begin{equation}
 C_{\min}(\tau)=\frac{51}{50}b^2e^{-1/\tau}
                  +o(e^{-1/\tau}).
 \label{eq:thermal-activation}
\end{equation}
The two groups have leading intensity ratio $9:8$. Reoptimizing the
exchange at each temperature changes only higher activation orders.
These four-cell lines identify the transitions that restore scattering
when the first excited triplet becomes populated.

\begin{figure*}[t]
\centering
\includegraphics[width=\textwidth]{thermal_evidence.pdf}
\caption{Thermal response at retuned interference zeros.
(a) Minimum leading positive-frequency coefficient in four- and six-cell
effective periodic chains. The doublet zero follows from cancellation of
the first-order operator. (b) Full microscopic two-cell weights
at $T/J_{\rm eff}=1.6$, normalized separately to their values at
$g=\beta/\Dloc=0.005$. Quartet and doublet slopes are $2.02549$ and
$4.00189$.
(c) Weights averaged over the uniform interval $q_0\pm\delta q$, relative
to their ideal central-momentum values, at $g=0.005$ and
$T/J_{\rm eff}=1.6$. The quartet and doublet central momenta are
$\pi$ and $3\pi/5$, respectively. Exchange parameters remain on their
central-momentum retuning trajectories. Both initial and final microscopic
states lie below $E_0+\Dloc/2$, while the partition function includes
the full spectrum. Lines connect computed points without smoothing.}
\label{fig:thermal-evidence}
\end{figure*}

\subsection{Microscopic spectra and the low-temperature crossover}

Full thermal traces on four- and six-cell effective chains give the
contrast in Fig.~\ref{fig:thermal-evidence}(a). At $\tau=1.6$, the quartet
minima are $0.00991078$ and $0.00968212$, respectively. For six cells the
latter is about 22 times its zero-temperature value, even after
temperature-dependent retuning.

The microscopic coupling dependence tests the predicted distinction
between second- and fourth-order weights.
We diagonalized all 1024 states of two coupled pentamers, including both
closing links and their corresponding effective exchange bonds.
The quartet root is Eq.~\eqref{eq:rank-exact-root}; the doublet uses
$q=3\pi/5$, $\alpha_*=1.09755143$, so both have $Q=\pi$.
Over $g=0.005,0.01,0.02$, the weights at $\tau=1.6$ have logarithmic
slopes $2.02549$ and $4.00189$ for the quartet and doublet, respectively
[Fig.~\ref{fig:thermal-evidence}(b)]. The quartet coefficient approaches
the effective prediction to within $1.18\%$ at the weakest coupling.
Appendix~\ref{app:exact-local-check} reports the temperature dependence
and finite-coupling errors.

The twenty-spin periodic model resolves the predicted four-cell activation.
We reconstructed its complete 256-state low band at
$g=0.000625,0.00125,0.0025$, keeping the exchange on the leading
zero-temperature trajectory $\kappa_0=1.28029750$.
At the weakest coupling the first
thermal gap is $1.000318J_{\rm eff}$; the two line-group centroids are
$2.000592J_{\rm eff}$ and $7.000309J_{\rm eff}$. Their unweighted
transition strengths, summed over the initial triplet and divided by
$\beta^2b^2$, are $0.536037$ and $0.482098$, within $0.74\%$ of
$27/50$ and $12/25$. Both approach their predictions as $g$ decreases
[Fig.~\ref{fig:four-cell-activation}(b)].

At finite coupling, a fourth-order ground-state background competes with
the exponentially small thermal contribution. At $\tau=0.1$, the
weakest-coupling weight exceeds the leading prediction by $25.9\%$;
ground-state transitions account for $21.0\%$ of the microscopic total.
This background scales as $\beta^{4.0015}$ over the three couplings and
is included in Fig.~\ref{fig:four-cell-activation}(a).
The response therefore crosses from thermal activation to a
ground-state background as temperature decreases at fixed $\beta$.

Higher-order splittings also produce quasielastic weight, amounting to
$1.68\%$ of the total at $\tau=1.6$.
The regular response with $\eta=0.05$ agrees with the effective
prediction within $0.23\%$ at the weakest coupling for
$\tau=0.2,0.4,1.6$. The resolved and unwindowed weights, together with
cutoff and solver checks, are specified in
Appendix~\ref{app:exact-local-check}.

\begin{figure*}[t]
\centering
\includegraphics[width=\textwidth]{four_cell_activation.pdf}
\caption{Direct four-cell microscopic check and low-temperature background.
(a) Resolved thermal weight for twenty physical spins with periodic
boundaries, $q=\pi$, fixed zero-temperature retuning, and
$\omega>0.05J_{\rm eff}$. Both initial and final states are in the
complete 256-state low band. The finite-coupling curves retain the
ground-state background. The effective curve is the full leading thermal
trace at the same retuning; the activation asymptote is
$(51/50)b^2\exp(-J_{\rm eff}/T)$.
(b) Strengths from the first thermally occupied triplet to the two predicted
final-state groups, divided by $\beta^2b^2(27/50)$ and
$\beta^2b^2(12/25)$, respectively. These strengths are summed over initial
magnetic components before thermal weighting. Lines connect calculated
points without smoothing.}
\label{fig:four-cell-activation}
\end{figure*}

\subsection{Momentum acceptance and window retuning}

A finite momentum interval samples nonzero form factors on either side
of the interference root, adding weight to the central response.
We averaged microscopic transition probabilities over a uniform window
around each root, initially keeping the central retuning fixed. At
$g=0.005,\tau=1.6$, a half-width $0.001\pi$ raises the quartet weight
by $0.094\%$, but the doublet weight by a factor of approximately
$2.93\times10^6$ [Fig.~\ref{fig:thermal-evidence}(c)].
The local doublet form factor has $\partial_qA=1.69998$ at its root,
whereas the quartet root at $q=\pi$ is stationary in momentum.
A fixed window therefore adds an ordinary form-factor background that
can overwhelm the doublet's very small fourth-order central signal.
At this fixed width, the doublet weight becomes nearly independent of
coupling, with logarithmic slope $0.00475$.

Optimizing the exchange for the entire window reduces, but does not
eliminate, this background. At the stationary quartet root, write
$x=q-\pi$, $A_{qq}=\partial_q^2 A(\alpha_*,\pi)$, and let $B$ be the
positive-frequency thermal weight of $D$. If $\kappa_0$ minimizes the
central response, the remaining source is thermally orthogonal to $D$.
The leading window-dependent dipole amplitude is
$\beta A'(\kappa-\kappa_0)+A_{qq}x^2/2$.
Retuning removes its mean, leaving the variance of $x^2$.
For a uniform half-width $\delta q=\rho\sqrt{\beta}$, at fixed finite
size and scaled temperature, this gives
\begin{align}
 \kappa_{\rm win}&=\kappa_0-\frac{A_{qq}\rho^2}{6A'}+o(1),\notag\\
 \frac{\overline W^{\rm opt}_{+,\eta}}{\beta^2}
 &\longrightarrow C_{\min}+\frac{A_{qq}^2B}{45}\rho^4.
 \label{eq:window-retuning}
\end{align}
Keeping $\kappa_0$ fixed gives $1/20$ in place of $1/45$ in the
background coefficient. Window retuning therefore leaves $4/9$ of this
leading background, set by the variance of the form factor.
In the ten-spin ring at $\tau=1.6$, the optimized microscopic
coefficients for $\rho=0.25,0.5,1$ agree with this limit within
$0.30\%$ at $g=0.00125$; deviations increase with coupling.

The chosen doublet root has nonzero $A_q=\partial_q A$, so its leading
window background is instead $A_q^2B(\delta q)^2/3$. Keeping it at
the fourth-order scale requires $\delta q=O(\beta^2)$, whereas the
stationary quartet root permits $O(\sqrt{\beta})$ at its quadratic
scale. The different tolerances follow from the slope and curvature
of these particular momentum zeros.
In the two-cell doublet, further exchange optimization cancels the
single low-band singlet--triplet amplitude to numerical precision.
Its weight for $\delta q=\rho\beta^2$ then approaches
$\beta^4 A_q^2B\rho^2/3$, within $0.055\%$ at the weakest tested coupling.
Thus the fourth power can also arise from momentum acceptance, even
when the central response has been suppressed further.
